Large language models (LLMs) are increasingly deployed as interacting agents that exchange 
information, negotiate decisions, and coordinate their actions over multiple turns. 
Such multi-agent systems (MAS) promise to combine information and capabilities distributed
across multiple agents, yet recent work has exposed substantial failures that arise 
specifically through interaction. Agents fail to proactively acquire information that 
is available only to others \citep{li2026systematic}, communicate actively without 
converting those exchanges into effective distributed reasoning \citep{zhang2026silobench}, 
and remain unable to coordinate effectively even when relevant information is made 
transparent \citep{yao2026talk}. Broader analyses similarly identify inter-agent failures 
as a distinct source of error beyond mistakes made by individual agents in isolation 
\citep{cemri2025multiagent,anthropic2026multiagent}. These findings raise a basic question: 
what capabilities are missing when individually capable agents fail to interact effectively?

Classical models of interactive decision making provide a principled way to make this question 
more precise. In an interactive POMDP, an agent maintains uncertainty not only over the 
physical environment, but also over models of other agents, and plans its own actions 
with respect to this interactive state \citep{gmytrasiewicz2005ipomdp}. Related work 
on opponent modeling similarly couples inference about another agent with planning a 
response to the inferred model \citep{huang2024hop}. Viewed through this lens, 
effective interaction requires at least two computations: \textbf{interactive-state inference}, 
which construct belief of other agents relevant information, and \textbf{interaction planning}, 
which uses that belief to evaluate what to do next. Importantly, these computations are 
coupled over time. An inferred state changes the value of a interaction action; 
the selected interaction then changes how other agents respond and therefore what 
evidence becomes available for subsequent inference. Established models of multi-agent 
decision making identify these two computations as fundamental components of interaction, 
motivating us to ask whether current LLM agents can reliably perform them together \citep{riemer2025tom}.

How can these capabilities be diagnosed and trained in a controlled way? 
Games provide a natural substrate because their information structure, utilities, legal 
actions, and strategic consequences can be specified and verified exactly. 
Recent work has accordingly used multi-agent games and self-play to train strategic 
or transferable reasoning in LLMs \citep{yuan2026marshal,liu2026spiral,lyu2026gift,he2026stratreasoner}. 
However, these environments are primarily designed around game-level performance rather than 
separate supervision of interactive-state inference and the planning that follows from it. 
Compact self-play games make scalable training possible but often abstract away the 
combination of latent information and heterogeneous incentives, whereas richer social games 
can entangle deduction, communication, deception, commitment and execution within the same 
trajectory \citep{light2023avalonbench}. To obtain precise supervision for the interaction 
structure above, we instead extend an existing commitment-based multi-party negotiation 
game \citep{benac2026benchmark}, which makes the environment becomes incomplete information and mixed incentives. 
Incomplete information makes relevant properties of other agents latent and therefore creates an inference problem, 
while mixed incentives make those inferred properties consequential for planning because the value of an interaction depends on what the other agents know, prefer, or are likely to do. 
Their combination captures an important class of real interactions in which agents act 
on behalf of different principals---for example, coordinating around private scheduling 
constraints \citep{zou2026calbench}, negotiating under asymmetric private objectives \citep{oneill2026cooperate}, 
or interacting repeatedly on behalf of distinct economic interests \citep{li2026emergent}.

Using this controlled setting, we separately diagnose failures in interactive-state inference 
and interaction planning, then construct strategic game slices that provide oracle 
supervision for both. We train on these slices and evaluate whether the resulting 
improvements generalize to unseen strategic configurations and transfer to external 
multi-agent environments. 

Our contributions are threefold:

\begin{itemize}
    \item \textbf{A controlled diagnosis of interactive reasoning in LLM-based MAS.} Grounded in classical interactive decision making, we study interactive-state inference and interaction planning in incomplete-information, mixed-incentive coordination, and separately measure failures in the two while retaining their sequential dependence.
    \item \textbf{A compact game construction for verifiable social-reasoning supervision.} We show how private preferences, consequential commitments, and information-revealing actions can instantiate a practically relevant interactive-decision problem while providing exact latent-state targets and counterfactual action values.
    \item \textbf{Targeted game supervision and transfer.} We construct strategic game slices around diagnosed interaction decisions, use their oracle structure to supervise inference and planning, and test whether the resulting capabilities generalize to unseen strategic configurations and external multi-agent environments.
\end{itemize}
